\section{Outlook}\label{sec:outlook}

\subsection{What the paper leaves}

Section~\ref{sec:generator} is settled: the generator is written out, the three
factors deliver the three terms of the target, and the whole shift dependence is
one $\cosh$ factor multiplying the off-diagonal data. Sections~\ref{sec:abelian}
and \ref{sec:endpoint} narrow the proposal rather than advancing it. The shift
direction is abelian; the two-parameter connection is flat wherever it exists;
the endpoint connection does not close on the form domain, and misses by exactly
the borderline. What survives is the atomic structure \eqref{eq:atomic}: a flat
connection with punctures at the prime periods, whose monodromy is currently
one-dimensional because the residues are scalars. Problems~\ref{prob:trace} and
\ref{prob:curvature} are the two smallest calculations that would move this.

\subsection{The decisive question}

The program's sharpest structural constraint is that the contraction whose norm
is Weil positivity is saturated, with a margin collapsing superexponentially in
$L$, so that any mechanism supplying a margin uniform in $L$ proves something
false \cite{Selection}. That statement is about the central form at $\omega=0$.
Proposition~\ref{prop:route} asks for something else: $\norm{V_{\omega_j,L_j}}\leq1$
along a sequence with $\omega_j\downarrow0$, which by \eqref{eq:storage} is a
statement about the cumulative defect
$\ip f{D_{\omega,L}f}=2\int_0^\omega Q_{s,L}[V_{s,L}f]\dd s$, not about its
integrand. The two differ: the cumulative defect can be positive on an input
where the instantaneous form is negative, and the program's own records contain
such a case at $L=\log7$, $\omega=10^{-11}$ \cite{Background}.

\begin{problem}[Does the flow formulation escape the criticality?]
\label{prob:criticality}
Decide whether a strict contraction $\norm{V_{\omega,L}}\leq1-\varepsilon$ at
fixed $\omega>0$, uniform in $L$, is consistent --- with the criterion, with the
explicit formula, and with the saturation at $\omega=0$. If it is, then the flow
formulation is the only formulation in this program in which a mechanism with a
margin is not immediately self-refuting, and that is the strongest available
argument for the whole direction. If it is not, the deformation picture inherits
the criticality unchanged and offers no relief from it.
\end{problem}

This should be settled before a gauge theory is chosen, not after. It is a
question about the arithmetic side alone.

\subsection{Order of work}

\begin{enumerate}
\item Problem~\ref{prob:criticality}. Cheapest to state, decisive for the
direction, and independent of any choice of theory.
\item Problem~\ref{prob:curvature}: the local factor at one atom, and what a
matrix-valued residue would have to satisfy.
\item Problem~\ref{prob:trace}: whether the endpoint flow closes anywhere.
\item Only then, a specific theory. Remark~\ref{rem:angle} says where to look
first --- a cusp, because the shift is an angle --- and
Section~\ref{sec:conditions} is the battery to run before computing anything
about a candidate. Condition~\ref{cond:marginal} already rejects the obvious
first attempt.
\end{enumerate}

\subsection{Scope}

Nothing above is proposed as a route to the Riemann hypothesis. Theorem
\ref{thm:kernel}, Lemma~\ref{lem:multiplicative}, Theorem~\ref{thm:abelian},
Proposition~\ref{prop:hyperbolic}, Proposition~\ref{prop:endpoint} and
Proposition~\ref{prop:flat} are written proofs; Proposition~\ref{prop:trace} is a
proof sketch, resting on a smoothing estimate quoted from \cite{Background} and
on the standard Sobolev trace threshold. Remark~\ref{rem:marginal} is an
observation about three exponents and explains nothing.
Section~\ref{sec:conditions}'s three conditions are proposals and are not part of
the companion investigation's index. No source is constructed, no positivity is
proved, no gauge theory is matched or excluded, and no statement here assumes the
Riemann hypothesis.

\subsection*{Preparation}

This working manuscript was produced with a large language model, Anthropic's
Claude (model \texttt{claude-opus-5}), working directly in the investigation
repository: the derivations and proofs of
Sections~\ref{sec:generator}--\ref{sec:endpoint}, the conditions of
Section~\ref{sec:conditions}, the check program of Appendix~\ref{app:checks}, and
the drafting throughout. The proposal the paper analyses is Edward Baker's. The
research notes accompanying the investigation name the model used for each. Every
statement is classified in the text as a written proof, a proof sketch, a
labelled numerical computation, or an assessment.
